% ============ 23 循环映射; Chern 类 (原书 pp.138-143) ============
\chapter{循环映射；Chern 类}\label{sec:23}
\index{循环映射}
\index{素循环}
\index{代数循环}
\index{Chow 环}
\index{Chern 类}
\index{Chern 特征}
\index{数值等价}
\index{Tate 猜想}
\index{半纯性}

我们希望把上同调类与簇上的代数循环联系起来.直接定义相当短，但遗憾的是，从它难以推出所需的性质.另一个使用 Chern 类的定义给出所有一般性质，但不那么显式.最好的方法是给出两个定义，并验证它们一致.

本节中，所有簇都在代数闭域 $k$ 上，$\Lambda = \mathbb{Z}/\ell^{n}\mathbb{Z}$，其中 $\ell \neq \operatorname{char}(k)$.我们也允许 $\Lambda$ 为 $\mathbb{Z}_{\ell}$ 或 $\mathbb{Q}_{\ell}$，尽管这需要对该论述作小的修改.我们设 $H^{*}(X) = \bigoplus_{r \geq 0} H^{2r}(X, \Lambda(r))$(记号如 \S 16 开头)---它在杯积下成为环.

作为参考，我们注意：给定簇与开子簇 $X \supset U \supset V$，得到平展上同调群的(三元组)正合序列
\[
\cdots \to H^{r}_{X \setminus U}(X) \to H^{r}_{X \setminus V}(X) \to H^{r}_{U \setminus V}(U) \to \cdots.
\]
当 $V$ 空时，这是空间对 $(X, U)$ 的正合序列.

\subsection*{预备}

设 $X$ 是 $k$ 上的非奇异簇.$X$ 上的素循环是不可约闭子簇.设 $C^{r}(X)$ 是由余维数 $r$ 的素循环生成的自由 Abel 群---它的元素称为 $X$ 上余维数 $r$ 的代数循环.于是 $C^{1}(X) = \mathcal{D}iv(X)$.我们设 $C^{*}(X) = \bigoplus_{r \geq 0} C^{r}(X)$.

关于两个代数循环有理等价的概念与两个代数循环的交积，我们请读者参考 Hartshorne 1977，附录 A.

$C^{*}(X)$ 除以有理等价的商 $CH^{*}(X)$ 在交积下成为环.它称为 Chow 环.

例如，$CH^{1}(X) = \operatorname{Pic}(X)$.回忆 $\operatorname{Pic}(X) = H^{1}(X_{\mathrm{et}}, \mathbb{G}_{m})$，因此 Kummer 序列的上同调序列给出同态
\[
\operatorname{Pic}(X) \to H^{2}(X_{\mathrm{et}}, \Lambda(1)).
\]
本节的目标是定义分次环的典范同态
\[
\mathrm{cl}^{*}_{X} \colon CH^{*}(X) \to H^{*}(X),
\]
使得 $\mathrm{cl}^{1}_{X} \colon CH^{1}(X) \to H^{2}(X_{\mathrm{et}}, \Lambda(1))$ 就是刚定义的映射.

\subsection*{循环映射的直接定义}

$X$ 又非奇异.当 $Z$ 非奇异时，设 $\mathrm{cl}_X(Z)$ 是 $1$ 在复合映射
\[
\Lambda = H^{0}(Z, \Lambda) \to H^{2r}_{Z}(X, \Lambda(r)) \to H^{2r}(X, \Lambda(r))
\]
(Gysin 映射)下的像.为把该定义推广到奇异素循环，需要如下引理.

\begin{lemma}\label{lem:23.1}[半纯性]
对 $X$ 中余维数 $c$ 的任何闭子簇 $Z$，对 $r < 2c$，$H^{r}_{Z}(X, \Lambda) = 0$.
\end{lemma}

\begin{proof}
当 $Z$ 非奇异时，$H^{s}_{Z}(X, \Lambda) \cong H^{s-2c}(Z, \Lambda(-c))$，对 $s - 2c < 0$ 它为零(见定理~\ref{thm:16.1}).我们对一般 $Z$ 的维数作归纳证明引理.若 $Z$ 的维数为 0，则它非奇异，从而论断成立.设 $Z$ 是余维数 $c$ 的闭子簇.它的奇异轨迹 $Y$ 是维数小于 $Z$ 的闭子簇，而三元组 $(X, X \setminus Y, X \setminus Z)$ 的正合序列是
\[
\cdots \to H^{r}_{Y}(X, \Lambda) \to H^{r}_{Z}(X, \Lambda) \to H^{r}_{Z \setminus Y}(X \setminus Y, \Lambda) \to \cdots.
\]
现在，因为 $Z \setminus Y$ 非奇异，对 $r < 2c$，$H^{r}_{Z \setminus Y}(X \setminus Y, \Lambda) = 0$；由归纳，对 $r < 2c + 2$，$H^{r}_{Y}(X, \Lambda) = 0$；这蕴含对 $r < 2c$，$H^{r}_{Z}(X, \Lambda) = 0$.
\end{proof}

现在设 $Z$ 是 $X$ 中余维数 $c$ 的素循环，$Y$ 是 $Z$ 的奇异轨迹.由三元组 $(X, X \setminus Y, X \setminus Z)$ 的正合序列得到(用引理)同构
\[
H^{2c}_{Z}(X, \Lambda) \cong H^{2c}_{Z \setminus Y}(X \setminus Y, \Lambda).
\]
我们定义 $\mathrm{cl}_X(Z)$ 为 $1$ 在复合映射
\[
\Lambda \cong H^{0}(Z \setminus Y, \Lambda(c)) \cong H^{2c}_{Z \setminus Y}(X \setminus Y, \Lambda(c)) \cong H^{2c}_{Z}(X, \Lambda(c)) \to H^{2c}(X, \Lambda(c))
\]
下的像.我们把它线性延拓为同态
\[
\mathrm{cl}^{r} \colon C^{r}(X) \to H^{2r}(X, \Lambda(r)).
\]
对 $r = 1$，这是 Kummer 序列定义的映射(见例~\ref{ex:7.9}a).

\subsection*{Chern 类}

现在取 $X$ 为射影非奇异簇.

设 $E$ 是 $k$ 上维数 $m + 1$ 的向量空间.我们把 $\mathbb{P}(E)$ 定义为 $E$ 中直线的射影空间.$E$ 的一组基的选取确定同构 $\mathbb{P}(E) \to \mathbb{P}^{m}$.更一般地，设 $E$ 是非奇异簇 $X$ 上秩 $m + 1$ 的 $\mathcal{O}_X$-模局部自由层(对 Zariski 拓扑).则(见 Hartshorne 1977, p.162)有与 $E$ 相联系的射影空间丛 $\mathbb{P}(E)$.它是带正则映射 $\varphi \colon \mathbb{P}(E) \to X$ 与典范可逆层 $\mathcal{O}(1)$ 的非奇异簇.纤维 $\varphi^{-1}(x) = \mathbb{P}(E_x)$，其中 $E_x$ 是 $E$ 在 $x$ 处的纤维($k$-向量空间，维数 $m + 1$).若 $E$ 在 $X$ 的 Zariski 开子集 $U$ 上自由，则同构 $\mathcal{O}^{m+1}_{U} \to E|_{U}$ 的选取确定同构 $\mathbb{P}^{m}_{U} \to \mathbb{P}(E|_{U})$.

\begin{theorem}\label{thm:23.2}
设 $E$ 是 $X_{\mathrm{zar}}$ 上秩 $m + 1$ 的局部自由层，$\varphi \colon \mathbb{P}(E) \to X$ 是相联系的射影丛.设 $\xi$ 是 $\mathcal{O}(1)$ 在 $H^{2}(\mathbb{P}(E), \Lambda(1))$ 中的类.则 $\varphi^{*}$ 使 $H^{*}(\mathbb{P}(E))$ 成为以 $1, \xi, \ldots, \xi^{m}$ 为基的自由 $H^{*}(X)$-模.
\end{theorem}

\begin{proof}
有分次环的同构
\[
\Lambda[T]/(T^{m+1}) \xrightarrow{\sim} H^{*}(\mathbb{P}^{m}),
\]
把 $T$ 送到超平面截的类 $\xi$；特别地，$H^{*}(\mathbb{P}^{m})$ 是以 $1, \ldots, \xi^{m}$ 为基的自由 $\Lambda$-模(见 \S 16).

若 $E$ 自由，同构 $\alpha \colon E \cong \mathcal{O}^{m+1}_{X}$ 的选取确定同构 $\mathbb{P}(E) \cong X \times \mathbb{P}^{m}_{k}$，而 K\"unneth 公式表明 $H^{*}(X \times \mathbb{P}^{m}_{k}) \cong H^{*}(X) \otimes H^{*}(\mathbb{P}^{m})$，它是 $H^{*}(X)$ 上以 $1, T, \ldots, T^{m}$ 为基的自由模.同构 $H^{*}(\mathbb{P}(E)) \cong H^{*}(X \times \mathbb{P}^{m}_{k})$ 是 $H^{*}(X)$-同构，且把 $\xi$ 送到 $T$.因此 $H^{*}(\mathbb{P}(E))$ 是 $H^{*}(X)$ 上以 $1, \xi, \ldots, \xi^{m}$ 为基的自由模.

其次，若 $U = X_0 \cup X_1$，且已知命题对 $X_0$，$X_1$ 与 $X_0 \cap X_1$ 成立，则 Mayer-Vietoris 序列(定理~\ref{thm:10.8})允许我们对 $U$ 证明它.由于已知命题对 $E$ 在其上平凡的任何 Zariski 开子集成立，这个论证允许我们逐步对 $X$ 得到它.
\end{proof}

该命题表明，当我们把 $H^{*}(\mathbb{P}(E))$ 看作 $H^{*}(X)$-模时，$1, \xi, \ldots, \xi^{m+1}$ 之间有线性关系，若把 $\xi^{m+1}$ 的系数归一化为 1，则它唯一.换言之，存在唯一元素 $\mathrm{ch}_r(E) \in H^{2r}(X, \Lambda(r))$，使得
\[
\begin{array}{c}
\sum_{r=0}^{m+1} \mathrm{ch}_r(E) \cdot \xi^{m+1-r} = 0\\
\mathrm{ch}_0(E) = 1.
\end{array}
\]
则 $\mathrm{ch}_r(E)$ 称为 $E$ 的第 $r$ 个 Chern 类，$\mathrm{ch}(E) = \sum \mathrm{ch}_r(E)$ 称为 $E$ 的全 Chern 类，而
\[
\mathrm{ch}_t(E) = 1 + \mathrm{ch}_1(E)t + \cdots
\]
称为 $E$ 的 Chern 多项式.

\begin{theorem}\label{thm:23.3}
Chern 类有如下性质：
\begin{enumerate}
  \item (函子性).若 $\varphi \colon Y \to X$ 是光滑簇的态射，$E$ 是 $X$ 上的向量丛，则 $\mathrm{ch}_r(\varphi^{*}(E)) = \varphi^{*}(\mathrm{ch}_r(E))$.
  \item (归一化).若 $E$ 是 $X$ 上的可逆层，则 $\mathrm{ch}_1(E)$ 是 $E$ 在 Kummer 序列给出的映射 $\operatorname{Pic}(X) \to H^{2}(X, \Lambda(1))$ 下的像.
  \item (可加性).若
  \[
  0 \to E' \to E \to E'' \to 0
  \]
  是光滑簇 $X$ 上向量丛的正合序列，则
  \[
  \mathrm{ch}_t(E) = \mathrm{ch}_t(E') \cdot \mathrm{ch}_t(E'').
  \]
\end{enumerate}
\end{theorem}

\begin{proof}
见 Grothendieck, A., \emph{La th\'eorie des classes de Chern}, Bull. Soc. Math. France 86 (1958), 137--158.
\end{proof}

上述参考文献也许看起来是时代错误，但 Grothendieck 证明的是抽象定理.他只假定存在
\begin{enumerate}
  \item[(i)] 从非奇异射影簇 $X$ 的范畴到分次反交换环的函子；
  \item[(ii)] 函子性同态 $\operatorname{Pic}(X) \to H^{2}(X)$;
  \item[(iii)] 对(光滑射影)簇的每个闭浸入 $Z \hookrightarrow X$，存在``Gysin 映射'' $H^{*}(Z) \to H^{*}(X)$,
\end{enumerate}
全都满足某些自然公理.然后他证明存在满足定理条件 (a, b, c) 的 Chern 类理论.

设 $K(X)$ 是 $X$ 上有限秩局部自由层的 Grothendieck 群.定理的 (c) 部分表明 $E \mapsto \mathrm{ch}(E)$ 穿过 $K(X)$.因为 $X$ 光滑，$K(X)$ 也是凝聚 $\mathcal{O}_X$-模的 Grothendieck 群(标准结果).这允许我们定义映射 $\psi \colon C^{*}(X) \to K(X)$：给定素循环 $Z$，用有限秩的自由 $\mathcal{O}_X$-模消解 $\mathcal{O}_Z$，
\[
0 \to E_n \to E_{n-1} \to \cdots \to E_0 \to \mathcal{O}_Z \to 0,
\]
并令 $\psi(Z) = \sum (-1)^{i}\mathrm{ch}(E_i)$(这里 $\mathcal{O}_Z$ 是 $Z$ 的结构层，看作支集在 $Z$ 上的 $\mathcal{O}_X$-模).总之，我们有映射
\[
C^{*}(X) \to K(X) \to H^{*}(X).
\]
为了得到分次群的同态，需要把 $K(X)$ 换成分次群.

$K(X)$ 上有滤过.定义 $K^{r}(X)$ 为支集在余维数 $\geq r$ 中的凝聚 $\mathcal{O}_X$-模生成的 $K(X)$ 的子群.群 $K^{r}$ 定义 $K(X)$ 的滤过.设
\[
GK^{*}(X) = \mathrm{gr}(K(X)) \overset{\mathrm{def}}{=} \bigoplus K^{r}(X)/K^{r+1}(X).
\]
则 $GK^{*}(X)$ 在乘积律
\[
[\mathcal{M}][\mathcal{N}] = \sum (-1)^{r}[\operatorname{Tor}^{\mathcal{O}}_r(\mathcal{M}, \mathcal{N})]
\]
下成为环.这里 $[\cdot]$ 表示 Grothendieck 群中的类.映射 $C^{*}(X) \to K(X)$ 显然与滤过相容，从而定义分次模的同态
\[
\psi \colon C^{*}(X) \to GK^{*}(X).
\]
设
\[
CH^{*}(X) = C^{*}(X)/(\text{有理等价}).
\]
它在交积下成为环.有理等价于零的循环在 $\psi$ 下映到零，因此 $\psi$ 定义映射 $\psi \colon CH^{*}(X) \to GK^{*}(X)$.Serre 对交积的描述表明这与交积相容.见 Hartshorne 1977，附录 A.

映射 $\mathrm{ch} \colon K^{*}(X) \to H^{*}(X)$ 诱导映射 $GK^{*}(X) \to H^{*}(X)$，但遗憾的是，它不是环同态：设 $H^{*}(X)_0$ 作为 Abel 群等于 $H^{*}(X)$，但赋予乘法结构
\[
x_r * x_s = \frac{(r+s-1)!}{(r-1)!(s-1)!} x_r x_s, \qquad x_r \in H^{r}(X), \ x_s \in H^{s}(X).
\]
则 $\mathrm{ch} \colon GK^{*}(X) \to H^{*}(X)_0$ 是同态.

当 $(2\dim X - 1)!$ 在 $\Lambda$ 中可逆时，映射
\[
x_r \mapsto \frac{x_r}{(-1)^{r-1}(r-1)!} \colon H^{*}(X)_0 \to H^{*}(X)
\]
是同构.复合映射
\[
CH^{*}(X) \to GK^{*}(X) \xrightarrow{\mathrm{ch}} H^{*}(X)_0 \to H^{*}(X),
\]
得到同态
\[
\mathrm{cl}_X \colon CH^{*}(X) \to H^{*}(X).
\]
它有下列性质：
\begin{enumerate}
  \item 它是分次环的同态(次数加倍)；特别地，交积映到杯积；
  \item 它关于 $X$ 函子性.
\end{enumerate}

\begin{theorem}\label{thm:23.4}
这个由 Chern 类定义的循环映射与直接定义的循环映射一致.
\end{theorem}

\begin{proof}
正确的证明相当长.
\end{proof}

\begin{exercise}\label{exc:23.5}
找出 EC VI 10.6 证明中的错误.
\end{exercise}

上面我们取了 $\Lambda = \mathbb{Z}/\ell^{n}\mathbb{Z}$.对不同的 $n$，循环映射相容，因此可以取逆向极限得到分次环的同态
\[
CH^{*}(X) \to \bigoplus H^{2r}(X, \mathbb{Z}_{\ell}(r)),
\]
只要 $\ell \geq 2\dim X$.与 $\mathbb{Q}$ 作张量积后，得到同态
\[
CH^{*}(X) \to \bigoplus H^{2r}(X, \mathbb{Q}_{\ell}(r))
\]
(不再需要 $\ell \geq 2\dim X$).

\begin{quote}
\textbf{注记} 本小节需要重写：来自 Carl Mautner 等人(UT Austin)的意见---讲义在这里有点混乱.你似乎漏掉了随后使用的 Chern 特征的定义.例如，在 23.3 之后你令 $\psi(Z) = \sum (-1)^{i}\mathrm{ch}(E_i)$，如果``ch''指 Chern 特征，这是有意义的且映到上同调(而不是你在上一段中说的 $K(X)$)，但你似乎忘了定义它.
\end{quote}

\subsection*{应用}

假定 $X$ 射影，维数为 $d$.群 $C^{d}(X)$ 由有限(形式)和 $\sum n_P P$ 组成，其中 $n_P \in \mathbb{Z}$，$P$ 是 $X$ 的闭点.次数映射
\[
\sum n_P P \mapsto \sum n_P \colon C^{d}(X) \to \mathbb{Z}
\]
穿过 $CH^{d}(X)$.因此，有配对
\[
Z, Y \mapsto (Z \cdot Y) \colon CH^{r}(X) \times CH^{d-r}(X) \to \mathbb{Z}.
\]
代数循环 $Z$ 称为数值等价于零，如果对所有余维数互补的代数循环 $Y$，$(Z \cdot Y) = 0$.我们设 $N^{r}(X)$ 为 $CH^{r}(X)$ 除以数值等价的商.现在有如下引人注目的定理.

\begin{theorem}\label{thm:23.6}
群 $N^{r}(X)$ 有限生成.
\end{theorem}

\begin{proof}
选取 $Y_1, \ldots, Y_n \in CH^{d-r}(X)$，使它们在 $H^{2d-2r}(X, \mathbb{Q}_{\ell}(d-r))$ 中的像构成由 $\mathrm{cl}(CH^{d-r}(X))$ 生成的 $\mathbb{Q}$-空间的一组基.同态
\[
Z \mapsto (Z \cdot Y_1, \ldots, Z \cdot Y_n) \colon CH^{r}(X) \to \mathbb{Z}^{n}
\]
的核恰好由数值等价于零的类组成，因此该同态定义单射 $N^{r}(X) \hookrightarrow \mathbb{Z}^{n}$.
\end{proof}

\begin{remark}\label{rem:23.7}
要证明 $N^{r}(X)$ 有限生成，只证明存在单射 $N^{r} \hookrightarrow H^{2r}(X, \mathbb{Q}_{\ell}(r))$ 不够.例如，存在单射 $\mathbb{Q} \hookrightarrow \mathbb{Q}_{\ell}$，但 $\mathbb{Q}$ 不是有限生成 Abel 群.
\end{remark}

设 $X$ 是可分闭域 $k$ 上的簇，$Z$ 是 $X$ 的不可约闭子簇.则因为 $X$ 与 $Z$ 由有限多个多项式定义，每个多项式只有有限多个系数，将存在定义在素数域上有限生成的域 $k_0$ 上的 $X$ 与 $Z$ 的模型.因此 $Z$ 的上同调类``定义在'' $k_0$ 上.Tate 猜想这可以用来刻画代数类.更精确地说，对 $X$ 在素数域上有限生成的域 $k_0$ 上的每个模型 $X_0$，设
\[
T^{r}(X_0/k_0) = H^{2r}(X_{\mathrm{et}}, \mathbb{Q}_{\ell}(r))^{\operatorname{Gal}(k^{\mathrm{sep}}_0/k_0)}.
\]
因为 $H^{2r}(X, \cdot) = H^{2r}(X_{0, k^{\mathrm{sep}}_0}, \cdot)$(见推论~\ref{cor:17.8})，这有意义.则 Tate 猜想
\[
T^{r}(X/k) \overset{\mathrm{def}}{=} \bigcup T^{r}(X_0/k_0)
\]
是 $H^{2r}(X, \mathbb{Q}_{\ell}(r))$ 的由代数循环的类生成的 $\mathbb{Q}_{\ell}$-子空间.虽然我们都希望猜想为真，但没有真正的理由相信它.对阿贝尔簇与 $r = 1$，它已被证明(Tate, Zarhin, Faltings).已知对 $\mathbb{C}$ 上的阿贝尔簇，Tate 猜想蕴含 Hodge 猜想(Piatetski-Shapiro, Deligne)，且已知对 $\mathbb{C}$ 上 CM 型阿贝尔簇的 Hodge 猜想蕴含对所有有限域上的阿贝尔簇的 Tate 猜想(所有 $r$)(Milne).然而，尽管 CM 型阿贝尔簇非常特殊---它们对应于模空间中零维的集---但似乎复分析学家/经典代数几何学家短期内不太可能证明它.
